% !TeX root = main.tex
\lecture{4}{Wed 28 Jan 2026 11:00}{Partial Differentiation IV: Total Time Derivative and Multivariate Taylor Series}

\section{Total Time Derivative}
Consider a fluid with a velocity field:
\[
    \underline{v}(x, y, z, t) = v_x(x, y, z, t) \hat{\imath} + v_y(x, y, z, t) \hat{\jmath} + v_z(x, y, z, t) \hat{k}
\]

We could follow a small piece of this fluid, and use Newton's Second Law, per unit volume. This gives mass per unit volume (i.e. density) times acceleration. Assuming a non-compressible fluid with constant density:

\[
    \underline{f} = \rho \frac{d \underline{v}}{dt}
\]

Where $\underline{f}$ is force per unit volume. Beginning by considering acceleration in only the x direction, there is two sources of time dependence. One from the fluid element itself moving, and one from the velocity of the piece changing with time.

This gives the total acceleration as:
\[
    \frac{d v_x}{dt} = \frac{\partial v_x}{\partial x} \frac{dx}{dt} + \frac{\partial v_x}{\partial y} \frac{dy}{dt} + \frac{\partial v_x}{\partial z} \frac{dz}{dt} + \frac{\partial v_x}{\partial t}
\]
\[
    = v_x \frac{\partial v_x}{\partial x} + v_y \frac{\partial v_x}{\partial y} + v_z \frac{\partial v_x}{\partial z} + \frac{\partial v_x}{\partial t}
\]

We can simplify this with nabla:
\[
    \frac{d v_x}{dt} = \left(v_x \frac{\partial}{\partial x} + v_y \frac{\partial}{\partial y} + v_z \frac{\partial}{\partial z} \right) v_x + \frac{\partial v_x}{\partial t} = \left(\underline{v} \cdot \underline{\nabla}\right) v_x + \frac{\partial v_x}{\partial t}
\]

The total time derivative or advective derivative is given by:
\[
    \frac{d \underline{u}}{dt} = \left(\underline{v} \cdot \nabla\right) \underline{u} + \frac{\partial \underline{u}}{\partial t}
\]

Where the first term gives the contribution from the element of fluid moving from the fluid flow pattern and the second term gives the contribution from the pattern itself changing.

The first term is called the advective term. A fluid element will accelerate even in a constant velocity field, because the velocity changes at it moves through the field.

The advective term is non-linear in the velocity field, which makes fluid mechanics difficult to solve.

\section{Multivariate Taylor Series}
For a single variable, we want to expand a function $f(x)$ around a point $x = a$ in a series of powers of $h$, where $h = x - a$:

\[
    f(a + h) = c_0 + c_1 h + c_2 h^2 + \cdots = \sum_{m=0}^{\infty} c_m h^m
\]

We find $c_m$ by differentiating m times wrt h and setting $h = 0$:
\[
    f^m(a) = m! c_m \implies c_m = \frac{f^m (a)}{m!}
\]

Hence:
\[
    f(a+h) = f(a) + hf^\prime(a) + \frac{h^2}{2!} f^{\prime\prime} (a) + \cdots = \sum_{m=0}^{\infty} \frac{h^m}{m!} f^m(a)
\]

Or in operator form:
\[
    f(a+h) = \exp \left[h \frac{\partial}{\partial x}\right] f(a)
\]

Noting that the exponential arises due to the Taylor series for the exponential function:
\[
    e^x = 1 + u + \frac{u^2}{2!} + \cdots = \sum_{m=0}^{\infty} \frac{u_m}{m!}
\]

\[
    \begin{aligned}
    f(a+h) &= \left. f(x) \right|_{x=a} \\
    &+ \left. h \frac{d}{dx} f(x) \right|_{x=a} \\
    &+ \left. \frac{h^2}{2!} \frac{d^2}{dx^2} f(x) \right|_{x=a} \\
    &\vdots \\
    &= \left. \left[ 1 + h \frac{d}{dx} + \frac{h^2}{2!} \frac{d^2}{dx^2} + \frac{h^3}{3!} \frac{d^3}{dx^3} \cdots \right] f(x) \right|_{x=a} \\
    &= \left. \exp \left[ h \frac{d}{dx} \right] f(x) \right|_{x=a}
\end{aligned}
\]


\subsection{Multiple Variables}
We want to expand a function $f(x, y)$ around the point $(x, y) = (a, b)$ as a power series.

This will take the form:
\[
    f(a+h, b+k) = c_{00} + c_{10} h + c_{01}k + c_{20} h^2 + c_{11}hk + c_{02}k^2 + \cdots \sum_{m=0}^{\infty} \sum_{n=0}^{\infty} c_{mn} h^m k^n
\]

We complete the same process to find coefficients, to find $c_{mn}$ we differentiate $m$ times wrt $k$, and $n$ times wrt $h$ and set $h = k = 0$ to get:

\[
    \frac{\partial^{m+n}}{\partial x^m \partial y^n} (a,b) = m!n!c_{mn}
\]

Hence:
\[
    f(a + h, b + k) = \sum_{m=0}^{\infty} \sum_{n=0}^{\infty} \frac{h^m k^n}{m!n!} \frac{\partial^{m+n} f}{\partial x^m \partial y^n} (a, b)
\]


Again in operator form this gives:
\[
    \exp \left[h \frac{\partial}{\partial x}\right] \exp \left[k \frac{\partial}{\partial y}\right] f(x, y) \biggr\rvert_{x=a, y=b}
\]

\subsection{Generalising to Many Variables}
Consider a function of $p$ variables $(x_1, x_2, \ldots, x_p)$:

\[
    f(a_1 + h_1, \ldots, a_p + h_p) = \sum_{n_1 = 0}^{\infty} \sum_{n_2 = 0}^{\infty} \cdots \sum_{n_p = 0}^{\infty} c_{n_1 n_ 2 \ldots n_p} h_1^{n_1} h_2^{n_2} \cdots h_p^{n_p}
\]

Again to find the coefficients we differentiate $n_1$ times wrt $h_1$, $n_2$ times wrt $h_2$ etc. We set $h_1 = h_2 = \cdots = h_p = 0$ to get:

\[
    n_1! n_2! \cdots n_p! c_{n_1 n_2 \ldots n_p} = \frac{\partial^{n_1 + n_2 + \cdots + n_p} f}{\partial x_1^{n_1} \partial x_2^{n_2} \cdots \partial x_p^{n_p}} (a_1, \ldots, a_p)
\]

Hence:
\[
    f(a_1 + h_1, \ldots, a_p + h_p) = \sum_{n_1 = 0}^{\infty} \cdots \sum_{n_p = 0}^{\infty} \frac{h_1^{n_1} \cdots h_p^{n_p}}{n_1! \cdots n_p!} \frac{\partial^{n_1 + \cdots + n_p} f}{\partial x_1^{n_1} \cdots \partial x_p^{n_p}} (a_1, \ldots, a_p)
\]

In compact vector notation, where $\underline{h} = (h_1, \ldots, h_p)$, this is:
\[
    f(\underline{a} + \underline{h}) = \exp \left[\sum_{i=1}^{p} h_i \frac{\partial}{\partial x_i}\right] f(\underline{a})
\]

\[
    f(\underline{a} + \underline{h}) = \exp \left[ \underline{h} \cdot \underline{\nabla} \right] f(\underline{a})
\]
